Each timed call follows the public certified compilation API, starting from
loaded syntax without reusing a proof or compiled program. Preparation
includes resolution and type checking. Certification/construction includes
semantic-plan preparation, admission, guard certification, and metadata
materialization. Lowering/validation includes executable lowering, metadata
emission, and final executable validation. Other cost covers dispatch and
intermediate-object cleanup. These are API boundaries; guard checking and
metadata construction can occur in more than one phase. Parsing, import
loading, serialization, portable loading, graph rewriting, and numerical
reference execution are outside timing.

The optimized compiler uses algebraic semantics and unchanged work budgets.
One warmup and ten measured rounds run serially, shuffling case order with
seed 20261026, on an Intel Core Ultra 7 155H. All original automatic rule
bodies, guards, and dialect bytes are preserved; imports are embedded.
The rank-three dot chains retain fixed source axis lists, enabling the
fixed-layout certificate; transpose permutations remain symbolic and use
the unary-chain certificate. Before timing, 60 independent dense-integer
workload checks and all 45 public check/compile/validate commands pass.
The driver also checks equality with public compilation and a portable
roundtrip for each of the 15 configurations.

Curves use medians with minimum--maximum whiskers; stacked phase bars use
means so their segments sum to mean total latency. One machine and ten
rounds characterize the tested workloads rather than an asymptotic bound
or production throughput. All 165 samples, exact sources, compiler hashes,
and validation logs remain in
\nolinkurl{evaluation/certified-scalability/}.
